% GENERATED FILE. Edit canonical lessons, then run npm run generate:books.
% canonical-source-hash: fnv1a64:f778915c2c77dd91
% canonical-lessons: HI-C460-virah, HI-C460-cimatna, HI-C460-akulana, HI-C460-suna, HI-C460-lalsa

\chapter{Missing Someone}
\label{ch:hi-missing-someone}
\hlchaptermodality{460}

\begin{chapteropening}
\textbf{By the end of this chapter:} \emph{I can say that I miss someone far away.}

Complete the last lesson of chapter 460: I can say that I miss someone far away.
\end{chapteropening}

\section[\texorpdfstring{\hi{विरह} (virah)}{विरह (virah)}]{\hi{विरह} (virah) — the pain of separation}

\label{lesson:HI-C460-virah}



\begin{quote}
Before the new one: say the Hindi for to kiss, then the Hindi for to soothe with soft words.
\end{quote}

\subsection*{You'll want to know: \hi{विरह}}
\textbf{\hi{विरह}} — \emph{virah} — \textquotedblleft{}the pain of separation\textquotedblright{}.

\begin{quote}
\hi{पति} \hi{के} \hi{विरह} \hi{में} \hi{पत्नी} \hi{बहुत} \hi{उदास} \hi{थी।}
\end{quote}

\begin{quote}
\emph{pati ke virah meṁ patnī bahut udās thī.}
\end{quote}

\begin{quote}
Parted from her husband, the wife was very sad.
\end{quote}

Masculine.

\begin{cousinweb}
\textbf{\hi{विरह}} (\emph{virah}) is a tatsama, borrowed from Sanskrit \emph{viraha}, \textquotedblleft{}parting, separation, esp. of lovers\textquotedblright{} (Monier-Williams). Songs and poems are full of it.
\end{cousinweb}

\begin{grammarlens}[title={What you've built}]
One more word for describing what you lived and felt.
\end{grammarlens}

\subsection*{Your turn}
\begin{itemize}
\raggedright
  \item \emph{Say it:} \emph{virah}
  \item \emph{Say it:} \emph{virah}, once more
  \item \emph{Say it:} \emph{pati ke virah meṁ patnī bahut udās thī.} — Parted from her husband, the wife was very sad.
  \item \emph{Recall:} say the Hindi for to kiss, then the Hindi for to soothe with soft words, then say \emph{virah} again
\end{itemize}

\begin{tcolorbox}[breakable,colback=teal!4,colframe=teal!35!black,title={Before you move on}]
What does \hi{विरह} mean? (\textquotedblleft{}The pain of separation\textquotedblright{}.) Say it once more.
\end{tcolorbox}
\section[\texorpdfstring{\hi{चिमटना} (cimaṭnā)}{चिमटना (cimaṭnā)}]{\hi{चिमटना} (cimaṭnā) — to cling tightly}

\label{lesson:HI-C460-cimatna}



\begin{quote}
Before the new one: say the Hindi for to soothe with soft words, then the Hindi for the pain of separation.
\end{quote}

\subsection*{You'll want to know: \hi{चिमटना}}
\textbf{\hi{चिमटना}} — \emph{cimaṭnā} — \textquotedblleft{}to cling tightly\textquotedblright{}.

\begin{quote}
\hi{स्टेशन} \hi{पर} \hi{बेटी} \hi{माँ} \hi{से} \hi{चिमट} \hi{गई।}
\end{quote}

\begin{quote}
\emph{sṭeśan par beṭī māṁ se cimaṭ gaī.}
\end{quote}

\begin{quote}
At the station the daughter clung tightly to her mother.
\end{quote}

Past (she …): \hi{वह} \hi{माँ} \hi{से} \hi{चिमट} \hi{गई}, \emph{vah māṁ se cimaṭ gaī}.

\begin{cousinweb}
\textbf{\hi{चिमटना}} (\emph{cimaṭnā}) is a tadbhava: CDIAL lists Hindi \emph{cimaṭnā}, \textquotedblleft{}to embrace closely\textquotedblright{}, under \emph{*cimb}, a form it reconstructs (CDIAL 4822).
\end{cousinweb}

\begin{grammarlens}[title={What you've built}]
One more word for describing what you lived and felt.
\end{grammarlens}

\subsection*{Your turn}
\begin{itemize}
\raggedright
  \item \emph{Say it:} \emph{cimaṭnā}
  \item \emph{Say it:} \emph{cimaṭnā}, once more
  \item \emph{Say it:} \emph{sṭeśan par beṭī māṁ se cimaṭ gaī.} — At the station the daughter clung tightly to her mother.
  \item \emph{Recall:} say the Hindi for to soothe with soft words, then the Hindi for the pain of separation, then say \emph{cimaṭnā} again
\end{itemize}

\begin{tcolorbox}[breakable,colback=teal!4,colframe=teal!35!black,title={Before you move on}]
What does \hi{चिमटना} mean? (\textquotedblleft{}To cling tightly\textquotedblright{}.) Say it once more.
\end{tcolorbox}
\section[\texorpdfstring{\hi{अकुलाना} (akulānā)}{अकुलाना (akulānā)}]{\hi{अकुलाना} (akulānā) — to grow restless; to pine}

\label{lesson:HI-C460-akulana}



\begin{quote}
Before the new one: say the Hindi for the pain of separation, then the Hindi for to cling tightly.
\end{quote}

\subsection*{You'll want to know: \hi{अकुलाना}}
\textbf{\hi{अकुलाना}} — \emph{akulānā} — \textquotedblleft{}to grow restless; to pine\textquotedblright{}.

\begin{quote}
\hi{बेटे} \hi{की} \hi{चिट्ठी} \hi{नहीं} \hi{आई}, \hi{तो} \hi{माँ} \hi{अकुला} \hi{गई।}
\end{quote}

\begin{quote}
\emph{beṭe kī ciṭṭhī nahīṃ āī, to māṁ akulā gaī.}
\end{quote}

\begin{quote}
When no letter came from her son, the mother grew restless with worry.
\end{quote}

Past (she …): \hi{वह} \hi{अकुला} \hi{गई}, \emph{vah akulā gaī}.

\begin{cousinweb}
\textbf{\hi{अकुलाना}} (\emph{akulānā}) is a tadbhava: CDIAL lists Hindi \emph{akulānā}, \textquotedblleft{}to be confused, be restless\textquotedblright{}, under Sanskrit \emph{ākula} (CDIAL 1012).
\end{cousinweb}

\begin{grammarlens}[title={What you've built}]
One more word for describing what you lived and felt.
\end{grammarlens}

\subsection*{Your turn}
\begin{itemize}
\raggedright
  \item \emph{Say it:} \emph{akulānā}
  \item \emph{Say it:} \emph{akulānā}, once more
  \item \emph{Say it:} \emph{beṭe kī ciṭṭhī nahīṃ āī, to māṁ akulā gaī.} — When no letter came from her son, the mother grew restless with worry.
  \item \emph{Recall:} say the Hindi for the pain of separation, then the Hindi for to cling tightly, then say \emph{akulānā} again
\end{itemize}

\begin{tcolorbox}[breakable,colback=teal!4,colframe=teal!35!black,title={Before you move on}]
What does \hi{अकुलाना} mean? (\textquotedblleft{}To grow restless; to pine\textquotedblright{}.) Say it once more.
\end{tcolorbox}
\section[\texorpdfstring{\hi{सूना} (sūnā)}{सूना (sūnā)}]{\hi{सूना} (sūnā) — desolate; empty}

\label{lesson:HI-C460-suna}



\begin{quote}
Before the new one: say the Hindi for to cling tightly, then the Hindi for to grow restless.
\end{quote}

\subsection*{You'll want to know: \hi{सूना}}
\textbf{\hi{सूना}} — \emph{sūnā} — \textquotedblleft{}desolate; empty\textquotedblright{}.

\begin{quote}
\hi{बच्चों} \hi{के} \hi{बिना} \hi{घर} \hi{सूना} \hi{लगता} \hi{है।}
\end{quote}

\begin{quote}
\emph{baccoṁ ke binā ghar sūnā lagtā hai.}
\end{quote}

\begin{quote}
The house feels empty without the children.
\end{quote}

\begin{cousinweb}
\textbf{\hi{सूना}} (\emph{sūnā}) is a tadbhava: CDIAL lists Hindi \emph{sūnā}, \textquotedblleft{}empty, hollow\textquotedblright{}, under Sanskrit \emph{śūnyá}, \textquotedblleft{}empty\textquotedblright{} (CDIAL 12567). It is said of places that feel empty because someone is missing.
\end{cousinweb}

\begin{grammarlens}[title={What you've built}]
One more word for describing what you lived and felt.
\end{grammarlens}

\subsection*{Your turn}
\begin{itemize}
\raggedright
  \item \emph{Say it:} \emph{sūnā}
  \item \emph{Say it:} \emph{sūnā}, once more
  \item \emph{Say it:} \emph{baccoṁ ke binā ghar sūnā lagtā hai.} — The house feels empty without the children.
  \item \emph{Recall:} say the Hindi for to cling tightly, then the Hindi for to grow restless, then say \emph{sūnā} again
\end{itemize}

\begin{tcolorbox}[breakable,colback=teal!4,colframe=teal!35!black,title={Before you move on}]
What does \hi{सूना} mean? (\textquotedblleft{}Desolate; empty\textquotedblright{}.) Say it once more.
\end{tcolorbox}
\section[\texorpdfstring{\hi{लालसा} (lālsā)}{लालसा (lālsā)}]{\hi{लालसा} (lālsā) — a longing; a craving}

\label{lesson:HI-C460-lalsa}



\begin{quote}
Before the new one: say the Hindi for to grow restless, then the Hindi for desolate.
\end{quote}

\subsection*{You'll want to know: \hi{लालसा}}
\textbf{\hi{लालसा}} — \emph{lālsā} — \textquotedblleft{}a longing; a craving\textquotedblright{}.

\begin{quote}
\hi{मीरा} \hi{को} \hi{अच्छी} \hi{नौकरी} \hi{की} \hi{लालसा} \hi{है।}
\end{quote}

\begin{quote}
\emph{Mīrā ko acchī naukrī kī lālsā hai.}
\end{quote}

\begin{quote}
Mira longs for a good job.
\end{quote}

Feminine.

\begin{cousinweb}
\textbf{\hi{लालसा}} (\emph{lālsā}) is a tatsama, borrowed from Sanskrit \emph{lālasā}, \textquotedblleft{}longing or ardent desire\textquotedblright{} (Monier-Williams).
\end{cousinweb}

\begin{grammarlens}[title={What you've built}]
One more word for describing what you lived and felt.
\end{grammarlens}

\subsection*{Your turn}
\begin{itemize}
\raggedright
  \item \emph{Say it:} \emph{lālsā}
  \item \emph{Say it:} \emph{lālsā}, once more
  \item \emph{Say it:} \emph{Mīrā ko acchī naukrī kī lālsā hai.} — Mira longs for a good job.
  \item \emph{Recall:} say the Hindi for to grow restless, then the Hindi for desolate, then say \emph{lālsā} again
\end{itemize}

\begin{tcolorbox}[breakable,colback=teal!4,colframe=teal!35!black,title={Before you move on}]
What does \hi{लालसा} mean? (\textquotedblleft{}A longing; a craving\textquotedblright{}.) Say it once more.
\end{tcolorbox}
